% 第三章：整体架构

\section{整体架构} \label{sec:design}

\subsection{分层与边界}

系统分三层：算法与模型、内核实现、实测载体（\autoref{fig:arch}）。算法只在第\ref{sec:theory}节定义一次；内核按它实现，实测用 QEMU 承载真实内核。证据链有两条，且不能互相替代：静态门禁证明"能编译、链接、符号一致"，内核实测证明"调度路径确实被调用、代价多大"。

\begin{figure}[H]
\centering
\begin{tikzpicture}[
  box/.style={draw,rounded corners=2pt,align=center,font=\small,minimum height=8mm,text width=3.6cm},
  side/.style={draw,dashed,rounded corners=2pt,align=center,font=\small,minimum height=8mm,text width=3.2cm,text=black!55},
  lab/.style={font=\small\bfseries,align=right},
  arr/.style={-{Latex[length=2mm]},thick},
  darr/.style={-{Latex[length=2mm]},thick,dashed,black!45}
]
\node[lab] at (-4.6,2.2) {算法/模型};
\node[box] (model) at (-1.6,2.2) {CRP 预测 + MSF/LCIF\\（第\ref{sec:theory}节）};
\node[lab] at (-4.6,0.6) {内核实现};
\node[box] (sched) at (-1.6,0.6) {调度类 \texttt{yat\_casched.c}\\选核 / 就绪池 / 历史表};
\node[box] (obs) at (2.6,0.6) {debugfs 计数器\\tracepoint};
\node[lab] at (-4.6,-1.0) {实测载体};
\node[box] (qemu) at (-1.6,-1.0) {QEMU + 自建 initramfs\\\texttt{yat\_bench}};
\node[box] (gates) at (2.6,-1.0) {\texttt{make verify-all}\\静态门禁};
\node[lab] at (-4.6,-2.6) {前期工作};
\node[side] (sim) at (-1.6,-2.6) {Java 模拟器\\（设计期校验，不作证据）};
\draw[arr] (model) -- (sched);
\draw[arr] (sched) -- (obs);
\draw[arr] (qemu) -- (sched);
\draw[arr] (gates) -- (sched);
\draw[darr] (model) -- (sim);
\end{tikzpicture}
\caption{系统分层：一份算法、一个内核实现、两条证据链（静态门禁 / 内核实测）；仿真器只作为前期工作的设计期校验}
\label{fig:arch}
\end{figure}

这条边界决定了本文的证据结构：第\ref{sec:test}节只报内核实测与静态门禁，不把仿真数据当作算法有效性的证据。

\subsection{内核侧组成}

调度类以内置方式编进内核（非模块），策略号为 8，配置项 \texttt{CONFIG\_SCHED\_CLASS\_YAT\_CASCHED}。集成点见\autoref{tab:kfiles}。

\begin{table}[H]
\centering
\caption{内核集成点}
\label{tab:kfiles}
\small
\begin{tabular}{@{}p{5.6cm}p{8.4cm}@{}}
\toprule
文件 & 改动职责 \\
\midrule
\texttt{kernel/sched/yat\_casched.c} & 调度类全部实现：就绪池、DAG/作业、历史表、加速表、选核、计数器 \\
\texttt{include/linux/sched/yat\_casched.h} & 数据结构与对外接口声明（29 个声明，符号检查门禁覆盖） \\
\texttt{include/trace/events/yat\_casched.h} & 三个 tracepoint：account / place / wakeup \\
\texttt{include/linux/sched.h} & \texttt{task\_struct} 增加 \texttt{yat\_casched} 实体 \\
\texttt{kernel/sched/core.c} & 策略校验、\texttt{sched\_class} 注册、初始化 \\
\texttt{kernel/sched/sched.h} & \texttt{yat\_casched\_policy()} 与 \texttt{rq} 扩展 \\
\texttt{kernel/sched/Makefile} & 编入 \texttt{yat\_casched.o} \\
\texttt{init/Kconfig} & 配置项与依赖 \\
\texttt{include/uapi/linux/sched.h} & \texttt{SCHED\_YAT\_CASCHED = 8} \\
\texttt{include/asm-generic/vmlinux.lds.h} & 链接段注册 \\
\bottomrule
\end{tabular}
\end{table}

\subsection{前期工作：仿真器}

项目前期实现过一个 Java 17 模拟器（\texttt{simulator/}，6361 行），用同一套缓存模型比较 WFD 与缓存感知启发式，并配有 CSV 导出、$t$ 区间统计与 116 项单测。它回答的是"在理想化模型下，某条放置规则会怎么表现"，而不是"在真实内核上会怎么表现"：其执行时间与命中判定由自建的缓存模型给出，不含抢占、迁移、TLB 与真实 LLC 竞争。本轮把它从证据链中移除，正文只在第\ref{sec:test}节用一小段交代结论与去处，代码与全部数据留在仓库文档中。

\subsection{复现入口}

\begin{verbatim}
make verify-all          # 静态门禁（含符号一致性、补丁包、文档链接）
make report              # 编译本文
make kernel-obj          # 只编调度器目标文件（秒级类型检查）
make kernel-build        # 全量编译，产出 arch/x86/boot/bzImage
scripts/run_kernel_repeats.sh --repeats 5     # 内核实测（3 组配置 x 5 次）
cd linux-6.8-main/yat_test_scripts && ./compare_cfs_vs_yat.sh \
   --tasks 8 --iters 2000000 --affinity none --sleep-us 200
\end{verbatim}

门禁刻意做成"快"：\texttt{make check} 查冲突标记与空文件，\texttt{make docs-check} 查 Markdown 链接，
\texttt{make name-check} 查项目名拼写，\texttt{make symbol-check} 查头文件声明与实现定义的一致性，
\texttt{make patch-check} 查补丁包 SHA256 与 hunk 数。这些检查都在秒级，可以在每次提交前跑完。
